%%%%%%%%%%  scan idx 115  |  folio 104  %%%%%%%%%%

We now consider localization and completion in the deformation theory
of schemes. We first fix notation. Let $X_0$ be a $k$-scheme of finite
type and let $x_0\in X_0$ be a fixed point. There are natural functors
\[
\underline M_{X_0}
\longrightarrow
\underline M_{\operatorname{Spec}(\mathcal O_{X_0,x_0})}
\longrightarrow
\underline M_{\operatorname{Spec}(\widehat{\mathcal O}_{X_0,x_0})},
\]

%%%%%%%%%%  scan idx 116  |  folio 105  %%%%%%%%%%

and hence functors
\[
\widehat{\underline M}_{X_0}
\longrightarrow
\widehat{\underline M}_{\operatorname{Spec}(\mathcal O_{X_0,x_0})}
\longrightarrow
\widehat{\underline M}_{\operatorname{Spec}
 (\widehat{\mathcal O}_{X_0,x_0})}.
\]
As observed in 4.1(c), $\widehat{\underline M}_{X_0}$ is canonically
equivalent to the category whose objects are pairs $(R,\mathfrak X)$,
where $R\in\operatorname{Ob}(\widehat{\underline C}_\Lambda)$ and
$\mathfrak X$ is a flat $R$-adic formal scheme equipped with a fixed
isomorphism $\mathfrak X_0\xrightarrow{\sim}X_0$.

When there is no danger of confusion, we denote the images of
$(R,\mathfrak X)$ under localization and completion by
$(R,\mathfrak X_{(x_0)})$ and
$(R,\mathfrak X_{(\widehat{x}_0)})$, respectively. Thus, if
\[
\mathfrak X=\varinjlim_\nu\mathfrak X_\nu,
\qquad
\mathfrak X_\nu
=\mathfrak X\times_{\operatorname{Spf}(R)}
 \operatorname{Spec}(R/\mathfrak m^{\nu+1}),
\]
then
\[
\mathfrak X_{(x_0)}
=\varinjlim_\nu
 \operatorname{Spec}(\mathcal O_{\mathfrak X_\nu,x_0})
\]
and
\[
\mathfrak X_{(\widehat{x}_0)}
=\varinjlim_\nu
 \operatorname{Spec}(\widehat{\mathcal O}_{\mathfrak X_\nu,x_0}).
\]
